\chapter*{Introduction}
\addcontentsline{toc}{chapter}{Introduction}

In the setting of a modern particle physics experiment, many cutting edge designs using various physical processes are utilized as detector systems. In the worlds largest circular particle accelerator experiment at CERN, there are 4 points of interest in which the oppositely accelerated proton beams are crossed at 13.4 TeV center of mass energy which creates a breeding ground for a large range of physics processes due to interactions from the sea of partons. At each point of interest (ALICE, ATLAS, LHCb and CMS), there is an intricate combined detector system set up to record the (aftermath) of these events.\\

Each one of these experiment points are designed to be sensitive to different specific particle physics processes and an overall classification of events. The detector subsystems that are first in line of the incoming particles from the interaction points are usually the tracker subsystems. Alongside systems like magnets, Cherenkov detectors, electromagnetic and hadronic calorimeters; the tracker systems are a vital part of a combined approach to analyze the plethora of particle physics processes. These tracker systems work by merely "noticing" and transferring the knowledge of the charged particles passing through in their proximity.\\

The LHCb detector when ordered from the closeness to the interaction point consists of several subsystems. Firstly, the VELO (Vertex Locator) system, which are silicon-strip detectors positioned close to the beam line to provide precise reconstruction of primary and secondary vertices. Then comes the Trigger Tracker station, a silicon microstrip detector that tracks particles before they enter the magnetic field from the massive dipole magnet which bends the paths of charged particles so that their momenta can be measured downstream. Following the magnet, the tracking stations which are composed of scintillating fibers, continue to track the charged particles as they exit the magnetic field, ensuring accurate momentum measurement. Further particle identification is done by the Ring-Imaging Cherenkov Detectors (RICH) which exploit the Cherenkov effect. Complementing these systems are the calorimeters including the muon calorimeters, which measure the energy of electrons, photons, hadrons and muons through the development of particle showers. Together, these subsystems enable LHCb to perform precise measurements and probe fundamental questions in particle physics.\cite{LHCb_detector}\\

In this laboratory course, we are going to study and try to understand the basics of the particular tracking system of scintillating fibers from the LHCb experiment upgrade in 2019/2020. 